\documentclass[a4paper,12pt]{article}

% ------------------------------------------------
% Packages
% ------------------------------------------------
\usepackage[utf8]{inputenc}
\usepackage[T1]{fontenc}
\usepackage[french]{babel}
\usepackage{geometry}
\usepackage{graphicx}
\usepackage{xcolor}
\usepackage{amsmath,amssymb}
\usepackage{tcolorbox}
\usepackage{fancyhdr}
\usepackage{enumitem}

\geometry{margin=2.5cm}

% ------------------------------------------------
% Informations personnalisables
% ------------------------------------------------
\newcommand{\university}{Votre Université}
\newcommand{\faculty}{Votre Faculté}
\newcommand{\module}{Statistiques avec R}
\newcommand{\level}{Licence / Master}
\newcommand{\teacher}{Dr. Aymen Ben Brik}
\newcommand{\doctype}{TP 0}
\newcommand{\docdate}{Année 2018--2019}

% ------------------------------------------------
% Header & Footer
% ------------------------------------------------
\pagestyle{fancy}
\fancyhf{}
\fancyhead[L]{\module}
\fancyhead[R]{\doctype}
\fancyfoot[C]{\thepage}

% ------------------------------------------------
% Page de garde personnalisée
% ------------------------------------------------
\newcommand{\makecustomtitle}{
\begin{center}

%\includegraphics[width=3cm]{logo.png} % Ajouter votre logo si besoin

\vspace{0.5cm}

{\Large \textbf{\university}}\\
{\faculty}

\vspace{0.5cm}

\rule{\linewidth}{0.5mm}

\vspace{0.3cm}

{\Huge \textbf{\doctype}}

\vspace{0.3cm}

{\Large \textbf{Rappels sur les statistiques avec R}}

\vspace{0.3cm}

\rule{\linewidth}{0.5mm}

\vspace{0.5cm}

\begin{tabular}{ll}
Module : & \module \\
Niveau : & \level \\
Enseignant : & \teacher \\
Année : & \docdate \\
\end{tabular}

\end{center}

\vspace{1cm}
}

% ------------------------------------------------
% Environnements personnalisés
% ------------------------------------------------
\newtcolorbox{objectif}{
colback=purple!5,
colframe=purple!60,
title=Objectif}

\newtcolorbox{exercicebox}{
colback=blue!5,
colframe=blue!60,
title=Exercice}

\newtcolorbox{questionbox}{
colback=green!5,
colframe=green!50,
title=Questions}

% ------------------------------------------------
% Document
% ------------------------------------------------
\begin{document}

\makecustomtitle

\begin{objectif}
Lors de ce TP, on va réviser les notions fondamentales de Statistiques avec le langage R : 
\begin{enumerate}
\item Traçage des courbes
\
\end{enumerate}
\end{objectif}

% ==================================================
\begin{exercicebox}
\textbf{Exercice 1}
\end{exercicebox}

\subsection*{Partie 1}

\begin{enumerate}
    \item Installer le package \textbf{Faraway}.
    \item Explorer le package \textbf{Faraway}.
    \item Charger le jeu de données nommé : \textbf{Coagulation}.
    \item Rédiger un petit résumé sur le jeu de données \textbf{coagulation}.
    \item Puisqu'on a 4 régimes d'alimentation différents, représenter graphiquement ``côte à côte'' les coagulations de chaque type de régime.
\end{enumerate}

\subsection*{Partie 2 : Applications}

\begin{enumerate}
    \item Rédiger un résumé sur le jeu de données \textbf{rats} issu du package \textbf{Faraway}.
    \item Faire la représentation graphique des données.
\end{enumerate}

% ==================================================
\begin{exercicebox}
\textbf{Exercice 2 : Les histogrammes}
\end{exercicebox}

\begin{enumerate}
    \item Représenter ces valeurs dans un histogramme :

    \begin{center}
    91, 49, 76, 112, 97, 42, 70, 100, 8, 112, 95, 90, 78, 62, 56, 94, 65, 58, 109, 70, 109, 91, 71, 76, 68, 62, 134, 57, 83, 66
    \end{center}

    \item Ajouter un titre à l'histogramme et une légende pour l'axe des abscisses.

    \item Représenter l'histogramme des fréquences (pas des effectifs).

    \item Changer la couleur de l'histogramme et représenter la courbe de densité dans le même graphique avec une autre couleur.
\end{enumerate}

% ==================================================
\begin{exercicebox}
\textbf{Exercice 3 : The Mauna Loa CO$_2$ DATA}
\end{exercicebox}

\begin{objectif}
\begin{itemize}
    \item Tracer une série temporelle.
    \item Utiliser R pour faire une régression linéaire.
    \item Tester la normalité des résidus.
    \item Tester la tendance de la droite de régression.
\end{itemize}
\end{objectif}

\begin{questionbox}
\begin{enumerate}
    \item Décrire les données contenues dans le jeu de données \textbf{CO2}.
    \item Représenter graphiquement les données.
    \item Déterminer les coefficients de la régression linéaire : CO$_2$ $\sim$ time(co2).
    \item Vérifier la normalité des résidus issus de la régression linéaire.
\end{enumerate}
\end{questionbox}

\end{document}